\documentclass[10pt,a4paper]{amsart}
\usepackage{amsmath,amssymb,amsthm,graphicx}
% Retained minimal preamble for the Stacks Project chapter more-morphisms.tex (More on Morphisms).
% Verbatim excerpt of the chapter preamble: Xy-pic setup (native preamble lines 14-18) and the native
% theorem declarations with the mathematical macros (native preamble lines 154-217). Copied mechanically;
% nothing added, removed or redefined. The wrapper class and the amsmath/amssymb/amsthm/graphicx base
% packages are supplied by the assembler, so the class/externaldocument part of the native preamble is omitted.
% For commutative diagrams we use Xy-pic
\usepackage[all]{xy}

% We use 2cell for 2-commutative diagrams.
\xyoption{2cell}
% Theorem environments.
%
\theoremstyle{plain}
\newtheorem{theorem}[subsection]{Theorem}
\newtheorem{proposition}[subsection]{Proposition}
\newtheorem{lemma}[subsection]{Lemma}

\theoremstyle{definition}
\newtheorem{definition}[subsection]{Definition}
\newtheorem{example}[subsection]{Example}
\newtheorem{exercise}[subsection]{Exercise}
\newtheorem{situation}[subsection]{Situation}

\theoremstyle{remark}
\newtheorem{remark}[subsection]{Remark}
\newtheorem{remarks}[subsection]{Remarks}

\numberwithin{equation}{subsection}

% Macros
%
\def\lim{\mathop{\mathrm{lim}}\nolimits}
\def\colim{\mathop{\mathrm{colim}}\nolimits}
\def\Spec{\mathop{\mathrm{Spec}}}
\def\Hom{\mathop{\mathrm{Hom}}\nolimits}
\def\Ext{\mathop{\mathrm{Ext}}\nolimits}
\def\SheafHom{\mathop{\mathcal{H}\!\mathit{om}}\nolimits}
\def\SheafExt{\mathop{\mathcal{E}\!\mathit{xt}}\nolimits}
\def\Sch{\mathit{Sch}}
\def\Mor{\mathop{\mathrm{Mor}}\nolimits}
\def\Ob{\mathop{\mathrm{Ob}}\nolimits}
\def\Sh{\mathop{\mathit{Sh}}\nolimits}
\def\NL{\mathop{N\!L}\nolimits}
\def\CH{\mathop{\mathrm{CH}}\nolimits}
\def\proetale{{pro\text{-}\acute{e}tale}}
\def\etale{{\acute{e}tale}}
\def\QCoh{\mathit{QCoh}}
\def\Ker{\mathop{\mathrm{Ker}}}
\def\Im{\mathop{\mathrm{Im}}}
\def\Coker{\mathop{\mathrm{Coker}}}
\def\Coim{\mathop{\mathrm{Coim}}}

% Boxtimes
%
\DeclareMathSymbol{\boxtimes}{\mathbin}{AMSa}{"02}

%
% Macros for moduli stacks/spaces
%
\def\QCohstack{\mathcal{QC}\!\mathit{oh}}
\def\Cohstack{\mathcal{C}\!\mathit{oh}}
\def\Spacesstack{\mathcal{S}\!\mathit{paces}}
\def\Quotfunctor{\mathrm{Quot}}
\def\Hilbfunctor{\mathrm{Hilb}}
\def\Curvesstack{\mathcal{C}\!\mathit{urves}}
\def\Polarizedstack{\mathcal{P}\!\mathit{olarized}}
\def\Complexesstack{\mathcal{C}\!\mathit{omplexes}}
% \Pic is the operator that assigns to X its picard group, usage \Pic(X)
% \Picardstack_{X/B} denotes the Picard stack of X over B
% \Picardfunctor_{X/B} denotes the Picard functor of X over B
\def\Pic{\mathop{\mathrm{Pic}}\nolimits}
\def\Picardstack{\mathcal{P}\!\mathit{ic}}
\def\Picardfunctor{\mathrm{Pic}}
\def\Deformationcategory{\mathcal{D}\!\mathit{ef}}
\newenvironment{reference}{}{}
\newenvironment{slogan}{}{}
\newenvironment{history}{}{}
\usepackage{cleveref}
\crefname{originallemma}{original lemma}{original lemmas}
\Crefname{originallemma}{Original lemma}{Original lemmas}
\makeatletter
\newlabel{algebra-lemma-regular-sequence-in-neighbourhood}{{061L}{source}{}{sourcepointer-1}{}}
\newlabel{algebra-lemma-regular-sequence-in-neighbourhood@cref}{{[originallemma][0][]061L}{[1][0][]source}{}{}{}}
\makeatother
\crefname{originallemma}{original lemma}{original lemmas}
\Crefname{originallemma}{Original lemma}{Original lemmas}
\makeatletter
\newlabel{divisors-lemma-Noetherian-scheme-regular-ideal}{{063I}{source}{}{sourcepointer-2}{}}
\newlabel{divisors-lemma-Noetherian-scheme-regular-ideal@cref}{{[originallemma][0][]063I}{[1][0][]source}{}{}{}}
\makeatother
\crefname{originallemma}{original lemma}{original lemmas}
\Crefname{originallemma}{Original lemma}{Original lemmas}
\makeatletter
\newlabel{dpa-lemma-perfect-map-ci}{{09PX}{source}{}{sourcepointer-3}{}}
\newlabel{dpa-lemma-perfect-map-ci@cref}{{[originallemma][0][]09PX}{[1][0][]source}{}{}{}}
\makeatother
\crefname{originallemma}{original lemma}{original lemmas}
\Crefname{originallemma}{Original lemma}{Original lemmas}
\makeatletter
\newlabel{more-algebra-lemma-noetherian-finite-all-equivalent}{{09CC}{source}{}{sourcepointer-4}{}}
\newlabel{more-algebra-lemma-noetherian-finite-all-equivalent@cref}{{[originallemma][0][]09CC}{[1][0][]source}{}{}{}}
\makeatother
\title{Avramov-type regularity: a perfect map into a target over a local complete intersection base forces both maps to be complete intersections}
\author{The Stacks Project Authors}
\date{stacks-project:more-morphisms.tex:local-snapshot-20261011; GNU FDL 1.2 or later}
\begin{document}
\maketitle
\noindent\textit{Editorial note.} Statement span is lines 18356-18374 of sources/stacks-more-morphisms/pinned-original-source.tex; proof span is lines 18376-18418, the complete author proof (41 body lines).\par\medskip
\noindent\textit{Editorial note.} Navigation stubs: 4 cited-but-not-included labels are supplied as non-load-bearing pointers with their pinned tags where the snapshot tag index has them.\par\medskip
\noindent\textit{Editorial note.} extra\_packages is empty on purpose: the retained preamble already loads xy with options and a second declaration clashes in the assembler wrapper.\par\medskip
\noindent\textit{Editorial note.} Text is retained verbatim, including the author's own typos and placeholders.\par\medskip
\noindent\textit{Editorial note.} This is the scheme-global form of Avramov's theorem that a perfect map with complete-intersection total-to-base map is itself a complete intersection; named by the author via the label.\par\medskip
\noindent\textit{Allowed standard prerequisites (editorial).}\begin{itemize}
\item Explicit prerequisite boundary. Equivalence of regular and Koszul-regular sequences in Noetherian rings (More on Algebra, lemma-noetherian-finite-all-equivalent), localization of regularity for ideals (Algebra, lemma-regular-sequence-in-neighbourhood; Divisors, lemma-Noetherian-scheme-regular-ideal), and the divided-powers criterion for perfect maps between regular immersions (Divided Power Algebra, lemma-perfect-map-ci).
\end{itemize}\medskip
\ifdefined\hypertarget\hypertarget{sourcepointer-1}{}\fi
\noindent\textit{Source reference.} Original lemma 061L: sources/stacks-more-morphisms/pinned-original-source.tex, source lines 18418--18418. Cited by the retained statement or proof but not part of the collected spans; supplied only as a navigation pointer. The argument does not depend on its content.\par\medskip
\ifdefined\hypertarget\hypertarget{sourcepointer-2}{}\fi
\noindent\textit{Source reference.} Original lemma 063I: sources/stacks-more-morphisms/pinned-original-source.tex, source lines 18418--18418. Cited by the retained statement or proof but not part of the collected spans; supplied only as a navigation pointer. The argument does not depend on its content.\par\medskip
\ifdefined\hypertarget\hypertarget{sourcepointer-3}{}\fi
\noindent\textit{Source reference.} Original lemma 09PX: sources/stacks-more-morphisms/pinned-original-source.tex, source lines 18418--18418. Cited by the retained statement or proof but not part of the collected spans; supplied only as a navigation pointer. The argument does not depend on its content.\par\medskip
\ifdefined\hypertarget\hypertarget{sourcepointer-4}{}\fi
\noindent\textit{Source reference.} Original lemma 09CC: sources/stacks-more-morphisms/pinned-original-source.tex, source lines 18418--18418. Cited by the retained statement or proof but not part of the collected spans; supplied only as a navigation pointer. The argument does not depend on its content.\par\medskip
\begin{lemma}
\label{lemma-lci-avramov}
Let
$$
\xymatrix{
X \ar[rr]_f \ar[rd] & & Y \ar[ld] \\
& S
}
$$
be a commutative diagram of morphisms of schemes. Assume
\begin{enumerate}
\item $S$ is locally Noetherian,
\item $Y \to S$ is locally of finite type,
\item $f : X \to Y$ is perfect,
\item $X \to S$ is a local complete intersection morphism.
\end{enumerate}
Then $X \to Y$ is a local complete intersection morphism
and $Y \to S$ is Koszul at $f(x)$ for all $x \in X$.
\end{lemma}

\begin{proof}
In the course of this proof all schemes will be locally Noetherian
and all rings will be Noetherian. We will use without further mention
that regular sequences and Koszul regular sequences agree in this
setting, see More on Algebra, Lemma
\ref{more-algebra-lemma-noetherian-finite-all-equivalent}.
Moreover, whether an ideal (resp.\ ideal sheaf) is regular
may be checked on local rings (resp.\ stalks), see
Algebra, Lemma \ref{algebra-lemma-regular-sequence-in-neighbourhood}
(resp.\ Divisors, Lemma \ref{divisors-lemma-Noetherian-scheme-regular-ideal})

\medskip\noindent
The question is local. Hence we may assume $S$, $X$, $Y$ are
affine. In this situation we may choose a commutative diagram
$$
\xymatrix{
\mathbf{A}^{n + m}_S \ar[d] & X \ar[l] \ar[d] \\
\mathbf{A}^n_S \ar[d] & Y \ar[l] \ar[ld] \\
S
}
$$
whose horizontal arrows are closed immersions. Let $x \in X$ be a
point and consider the corresponding commutative diagram of local
rings
$$
\xymatrix{
J \ar[r] &
\mathcal{O}_{\mathbf{A}^{n + m}_S, x} \ar[r] &
\mathcal{O}_{X, x} \\
I \ar[r] \ar[u] &
\mathcal{O}_{\mathbf{A}^n_S, f(x)} \ar[r] \ar[u] &
\mathcal{O}_{Y, f(x)} \ar[u]
}
$$
where $J$ and $I$ are the kernels of the horizontal arrows.
Since $X \to S$ is a local complete intersection morphism, the
ideal $J$ is generated by a regular sequence. Since $X \to Y$ is
perfect the ring $\mathcal{O}_{X, x}$ has finite tor dimension over
$\mathcal{O}_{Y, f(x)}$. Hence we may apply
Divided Power Algebra, Lemma \ref{dpa-lemma-perfect-map-ci}
to conclude that $I$ and $J/I$ are generated by regular sequences.
By our initial remarks, this finishes the proof.
\end{proof}

\end{document}
